\begin{helper}
Fix a history cell with recorded coordinates \(R_0\), terminal coordinates
\(U\), and a nodup reveal order \(o\) with \(\operatorname{toFinset}(o)=U\).
Assume both \(R_0\) and \(U\) are contained in the terminal coordinate
universe \(\noderef{TwoBiteTerminalCoordinateUniverse}(m)\), and assume
terminal coordinates are disjoint from the recorded coordinates.
Assume the branch/transcript prefix-safety hypothesis along \(o\): if two
samples agree on the embedding, all recorded coordinates, and all terminal
coordinates before \(e\) in \(o\), then at \(e\) they have the same
staged-open, touched, same-color-closed, and pre-terminal-recorded status.

Fix support labels \(B,J\), realized pairs
\(\mathsf{Realized}(B,J,\beta,\tau)\), transcript triples
\(\tau=(P_\tau,A_\tau,Z_\tau)\), complete blue traces
\(\operatorname{blueTrace}(\beta)\), fixed red supports
\(\operatorname{redSupport}(\beta,J)\), and complete-assignment
compatibility predicates \(\mathsf{Comp}(A;B,J,\beta,\tau)\).  Assume every
realized pair is an allowed transcript label and has the exported geometry:
the red support is contained in \(U\) and red-only; \(P_\tau,A_\tau\subseteq U\);
\(P_\tau\cap A_\tau=\emptyset\);
\[
B\cup\operatorname{redSupport}(\beta,J)\subseteq P_\tau,\qquad
Z_\tau\subseteq A_\tau,
\]
and \(P_\tau,A_\tau\) encode the complete terminal blue trace.  Assume also
the forward compatibility consequences: a compatible complete assignment
contains \(P_\tau\), avoids \(A_\tau\setminus Z_\tau\), and agrees with
\(\operatorname{blueTrace}(\beta)\) on all blue terminal coordinates.

Assume moreover that the fixed-\((B,J)\) scan data have been supplied
explicitly.  Thus there is a partial, transcript-valued candidate scan, a
witness terminal truth table \(W_{\beta,\tau}\), exact read sets
\(D_{\beta,\tau}(e)\), and four scan tests: staged-open, touched,
same-color-closed, and pre-terminal-recorded.  For every realized
\((\beta,\tau)\), the witness lies in \(U\), contains every coordinate of
\(P_\tau\), and contains no coordinate of \(A_\tau\).  Every read coordinate
is classified as fixed present, residual present, selected absent, or
unselected absent; each read set lies in the earlier prefix of \(o\); agreement
with the witness on the read set preserves all four tests; and agreement of
all four tests at every terminal coordinate forces the supplied transcript
scan output \(\tau\).  No total assignment-to-branch map is required; the
branch label used in each scan calculation is the already-realized label
\(\beta\).

Finally assume the overlap/pruning interface needed for residual
single-ownership: prefix sets lie in \(U\) and make the four tests prefix
stable; complete blue traces determine realized branch labels;
selected-present siblings are selected coordinates; the prefix fixed support
and open-candidate geometry, sibling data, support exhaustion, and the
first-mismatch alternative all hold for the fixed \(B,J\).  This is a
realized-only interface: overlap pruning uses only compatibility of the
realized cells under consideration, rather than a total branch assignment for
every terminal truth table.

Then the supplied candidate scan witnesses the existence of a deterministic
transcript scan
\[
\operatorname{scan}_{B,J}:2^U\to\mathcal P\cup\{\bot\}
\]
with the three fixed-\((B,J)\) interface clauses needed by
\(\noderef{FixedSetHistoryCellRedResidualDeterministicScanOutputs}\).
First, residual delete-and-reinsert stability holds: if
\[
P_\tau\setminus(B\cup\operatorname{redSupport}(\beta,J))
  \subseteq A_{\rm res},\qquad
A_{\rm res}\cap(A_\tau\setminus Z_\tau)=\emptyset,
\]
then
\[
\operatorname{scan}_{B,J}\bigl((A_{\rm res}\setminus Z_\tau)\cup
  (B\cup\operatorname{redSupport}(\beta,J))\bigr)=\tau .
\]
Second, scan recognition holds under the full deterministic scan data for a
realized \((\beta,\tau)\): if a terminal assignment contains every forced
present coordinate in \(P_\tau\), contains no forced absent coordinate in
\(A_\tau\), and has the complete blue trace of \(\beta\), then the scan
returns \(\tau\).  The complete blue trace is therefore only one part of the
recognition input; it is not asserted to determine \(\tau\) by itself.
Third, if one residual assignment lies in the residual cylinders of two
realized pairs and the two delete-and-reinsert complete assignments are
compatible with those pairs, then the two branch/transcript owners are equal.
\end{helper}
\begin{proof}
Use the supplied partial transcript scan, whose data are generated by the same
blue-first/red-second reveal order used in the paper.  The scan reads the
terminal coordinates in the order \(o\),
first through the blue part and then through the red part.  At a coordinate
\(e\), the only information used by the scan is the earlier terminal
prefix: staged-openness, touchedness, same-color closure, and
pre-terminal-recordedness at \(e\).  The hypothesis on prefix safety says
that these four tests are stable under replacing the comparison sample by
any other terminal truth table that agrees on the earlier prefix and on the
recorded coordinates.

The comparison truth table is made into a genuine two-bite sample by
\(\noderef{FixedSetHistoryCellTerminalAssignmentExtension}\).  Its recorded
coordinates agree with the realized witness, its embedding is unchanged,
and its terminal coordinates are exactly the prescribed terminal assignment.
Thus the prefix-safety hypothesis applies to the comparison samples used by
the scan, not merely to abstract Boolean strings.

For a realized transcript \(\tau\), let \(D_\tau(e)\) be the exact set of
earlier terminal coordinates whose values are read when the scan evaluates
the four tests at \(e\).  The construction of
\(\noderef{FixedSetHistoryCellRISIRevealOrderScanConstruction}\) supplies
the generated read-set interface: each \(D_\tau(e)\) lies in the earlier
prefix of \(o\), every read coordinate is classified by the transcript as
fixed present, variable present, selected absent, or unselected absent, and
agreement on \(D_\tau(e)\) preserves the four tests at \(e\).  In the
present notation this means
\[
D_\tau(e)\subseteq P_\tau\cup A_\tau .
\]
The deterministic-output clause says that once the four test outcomes agree
at every terminal coordinate for a realized \((\beta,\tau)\), the ordered
transcript scan returns \(\tau\).  In the formal proof the generic scan lemma
is instantiated locally with the constant branch \(\beta\), which is legal
because the calculation is already inside a realized \(\beta\)-cell; no
global inhabitant of the branch-label type is chosen.

Now take a residual assignment \(A_{\rm res}\) satisfying the displayed
cylinder conditions and form
\[
A^\ast=(A_{\rm res}\setminus Z_\tau)\cup
  (B\cup\operatorname{redSupport}(\beta,J)).
\]
Every present read coordinate is present in \(A^\ast\): fixed present
coordinates are reinserted through
\(B\cup\operatorname{redSupport}(\beta,J)\), and variable present
coordinates lie in \(A_{\rm res}\) by the residual containment.  Every
absent read coordinate is absent in \(A^\ast\): unselected absences are
excluded by the residual avoidance, and selected absences are deleted by
subtracting \(Z_\tau\).  Hence the realized witness and \(A^\ast\) agree on
each \(D_\tau(e)\).  Applying
\(\noderef{FixedSetHistoryCellRISIResidualScanStability}\) gives the first
scan clause,
\[
\operatorname{scan}_{B,J}(A^\ast)=\tau .
\]

For scan recognition, let \(A\subseteq U\) agree with the full truth table
recorded by a realized \((\beta,\tau)\): every coordinate of \(P_\tau\) is
present in \(A\), every coordinate of \(A_\tau\) is absent from \(A\), and
the blue terminal coordinates of \(A\) are exactly
\(\operatorname{blueTrace}(\beta)\).  The blue-first part of the reveal order
fixes the branch from the complete terminal blue truth table.  The
red-second part then evaluates the same generated read sets \(D_\tau(e)\)
against the comparison sample for \(A\).  Because
\(\noderef{FixedSetHistoryCellRISIRevealOrderScanConstruction}\) classifies
every read coordinate as forced present or forced absent for \(\tau\), the
comparison sample and the realized witness agree on every read set
\(D_\tau(e)\).  Exact-read stability therefore gives the same ordered
sequence of test outcomes, and the deterministic-output clause of the same
node forces the scan output on \(A\) to be \(\tau\).  This is the formal
version of the paper's statement that after all of \(G_B\) has been revealed,
the set from which the red choices are drawn is completely determined, with
the red-second scan data included explicitly.

Finally suppose a residual assignment reconstructs compatibly to two
realized pairs.  The forward compatibility consequences convert each
compatibility certificate into the corresponding residual-cylinder
containment and avoidance conditions.  If the two owners were different,
the first-mismatch certificate from
\(\noderef{FixedSetHistoryCellFixedBJResidualOverlapConsistency}\), using
the prefix invariant of
\(\noderef{FixedSetHistoryCellRedPrefixScanInvariant}\), would produce a
selected-present coordinate in one reconstructed assignment before it is
legally released in the other.  The pruning conclusion then rules out the
very compatibility certificate for the realized cell that produced the
selected-present sibling.  This contradiction proves that the two realized
owners are equal.
\end{proof}
